\begin{abstract}
Syniscopy is a shared-scene Fisher-information diagnostic for comparing configured nanoparticle microscopes under explicit forward-model assumptions. The unit of comparison is a \emph{microscope}: a chosen contrast modality, which fixes the physics backend, together with that instrument's own complete optical, detector, and backend configuration. Two microscopes may select the same modality yet differ entirely in their parameters, and instrument quantities such as numerical aperture or pupil sampling belong to a microscope rather than to a modality. The method maps multiple configured microscopes onto a common rigid-particle scene parameterization, including lateral position, axial position, and local \(\mathrm{SE}(3)\) pose coordinates, then computes Cramer--Rao diagnostics per microscope: lateral derivatives use the exact spectral gradient of the sampled band-limited contrast image, while axial and orientation derivatives remain explicit symmetric rerenders. The workflow emits status-gated artifacts for microscope ranking, detected-quanta normalization, algebraic fusion, and supervision-mask policy checks. Comparing bare modalities is recovered as the special case of fixing one instrument and varying only its contrast modality. The central claim is not a universal benchmark ordering; it is that comparisons are interpretable only after the scene, instrument configuration, measurement domain, noise model, and derivative contract are fixed. We demonstrate the workflow on configured optical and electron profile cards, native-regime source-provenance audits, and deterministic supervision-policy cases. Headline paper artifacts are tied to explicit sampling, convergence, and validation-status files so lateral undersampling/support failures and axial/orientation derivative gates remain visible rather than being hidden in prose.
\end{abstract}
